\section{Künstliche Architekturen}
\slide{Begrifflichkeiten}{
    \begin{itemize}
        \item \b{Künstliche Intelligenz (KI):} Forschungsfeld, das menschliche Wahrnehmungs- und Denkprozesse als Berechnungsverfahren versteht und versucht, diese in Computern zu implementieren oder intelligente Systeme zu bauen.
        \item \b{Kaskadielle Architektur:} Eine Programmstruktur, bei der alle Module gleichzeitig aktiv sind und Daten seriell von einem Modul zum nächsten wandern (ähnlich einem Wasserfall).
        \item \b{Hebb-Regel:} Lernregel in neuronalen Netzen: "What fires together, wires together." Wenn zwei Knoten gleichzeitig aktiviert sind, verstärkt sich ihre Verbindung (Gewichtung der Kante).
    \end{itemize}
}
\slide{Philosophische Konzepte \& Kritik}{
    \begin{itemize}
        \item \b{Turing-Test:}
        \begin{itemize}
            \item Ablauf: Ein Mensch kommuniziert (z. B. per Tastatur) mit einem unsichtbaren Partner (Mensch oder Computer).
            \item Ziel: Kann der Mensch nicht unterscheiden, ob er mit einem Computer oder einem Menschen spricht, gilt der Computer als intelligent.
        \end{itemize}
        \item \b{Chinese Room (Searle):}
        \begin{itemize}
            \item Gedankenexperiment: Jemand sitzt in einem Zimmer und sortiert chinesische Zeichen nach Regeln, ohne Chinesisch zu verstehen. Für Außenstehende sieht der Output perfekt aus.
            \item Kritik an KI: Bloße Symbolverarbeitung (Syntax) bedeutet kein inhaltliches Verstehen (Semantik) oder Bewusstsein.
        \end{itemize}
        \item \b{KI-Thesen:}
        \begin{itemize}
            \item Schwache KI: Computer simulieren intelligentes Verhalten nur; sie verstehen nicht wirklich, was sie tun (Symbolverarbeitung).
            \item Starke KI: Computer können prinzipiell genauso denken, verstehen und handeln wie Menschen (Geist = Programm).
        \end{itemize}
    \end{itemize}
}
\slide{Kognitive Architekturen}{
    \begin{itemize}
        \item \b{Blackboard-Architektur:}
        \begin{itemize}
            \item Mehrere unabhängige Module arbeiten parallel und ohne zentrale Kontrolle.
            \item Kommunikation läuft über einen gemeinsamen Datenspeicher, das "Blackboard", auf das alle Module zugreifen können.
        \end{itemize}
        \item \b{ACT-R (Hybrides Modell):}
        \begin{itemize}
            \item Kombiniert symbolische (Regeln) und subsymbolische (Aktivierung) Ansätze.
            \item Langzeitspeicher: Besteht aus Prozeduralem Gedächtnis (Regelspeicher für Fertigkeiten/Abläufe) und Deklarativem Gedächtnis (Faktenwissen in semantischen Netzwerken).
        \end{itemize}
    \end{itemize}
}
\slide{Künstliche Neuronale Netze (KNN)}{
    \begin{itemize}
        \item Funktionsweise: Orientiert an biologischen Neuronen; Verarbeitung durch Aktivationsausbreitung über gewichtete Kanten (parallele Verarbeitung).
        \item Lokalisierte Repräsentation: Ein einzelner Knoten steht für eine bestimmte Bedeutung (z. B. ein Knoten = Konzept "Hund").
        \item \b{Distribuierte Repräsentation:} Ein Konzept wird durch das Aktivierungsmuster vieler Knoten dargestellt; einzelne Knoten haben keine isolierte Bedeutung.
    \end{itemize}
}
